\bookchapter{บทที่ 23 ภาคผนวก: คู่มือเริ่มต้น GitHub — จากที่เก็บโค้ดสู่แพลตฟอร์มการทำงานร่วมกัน}

\chapterintro{บทที่แล้วเราคุยเรื่อง “การคิดแบบโปรแกรมเมอร์” กันไปแล้ว แต่จะเอาแนวคิดนั้นมาลงมือทำจริงและร่วมงานกับคนอื่น ก็ข้ามชื่อหนึ่งไปไม่ได้ นั่นคือ GitHub ไม่ว่าคุณจะเป็นนักศึกษาคอมพิวเตอร์กำลังเตรียมแข่งขัน อยากเก็บผลงานโค้ดของตัวเอง หรือแค่อยากแอบดูว่านักพัฒนาทั่วโลกกำลังเล่นอะไรกันอยู่ GitHub ก็เป็นแพลตฟอร์มที่คุณต้องรู้จัก ในหัวข้อนี้เราจะเริ่มจากคอนเซปต์พื้นฐานที่สุด เพื่อช่วยให้คุณเข้าใจมันได้ชัดเจน}

\booksection{หนึ่ง GitHub แท้จริงแล้วคืออะไร?}

พูดง่าย ๆ ก็คือ GitHub เป็น\textbf{แพลตฟอร์มจัดเก็บโค้ดและทำงานร่วมกัน}ที่สร้างบนพื้นฐานของ Git เปิดตัวปี 2008 และถูกไมโครซอฟต์ซื้อกิจการในปี 2018 ปัจจุบันยังคงเป็นที่เก็บโค้ดโอเพนซอร์สที่ใหญ่ที่สุดในโลก ดูแลรีโปอยู่หลายร้อยล้านแห่ง ครอบคลุมภาษาโปรแกรมกระแสหลักแทบทุกภาษา

ถ้าจะใช้การเปรียบเทียบเพื่อความเข้าใจ ก็คือ:\textbf{ถ้า Git คือสมุดบันทึกบนคอมของคุณ GitHub ก็คือแพลตฟอร์มบล็อกสาธารณะ}— คุณเผยแพร่ผลงานของตัวเอง ติดตามโปรเจกต์ของคนอื่น ร่วมพูดคุยและทำงานร่วมกันได้ มันไม่ใช่แค่ “คลังหลังบ้านของโปรแกรมเมอร์” แต่คล้าย “เมืองเทคโนโลยี” ที่นักพัฒนาทั่วโลกร่วมกันสร้าง — โปรเจกต์ถือกำเนิด พัฒนาต่อ ถูกถกเถียง และถูกนำไปใช้ซ้ำที่นี่

สำหรับคนทั่วไป GitHub ก็ไม่ได้ไกลตัวอย่างที่คิด คนที่ไม่ใช่โปรแกรมเมอร์จำนวนมากใช้มันเก็บโน้ต อีบุ๊ก บล็อกส่วนตัว หรือแม้แต่สำรองไฟล์ตั้งค่าต่าง ๆ แค่คุณมีความต้องการนิดหน่อยที่จะ “จัดระเบียบของตัวเองให้เข้าที่แล้วแชร์ออกไป” GitHub ก็คุ้มค่าที่จะรู้จักแล้ว

\booksection{สอง ความต่างระหว่าง Git กับ GitHub (คอนเซปต์สำคัญ!)}

หลายคนสับสนสองคำนี้ตั้งแต่ครั้งแรกที่เจอ มาแยกให้ชัดกันเลย:

\begin{center}
\begin{tabularx}{\linewidth}{|X|X|X|}
\hline
\textbf{หัวข้อเปรียบเทียบ} & \textbf{Git} & \textbf{GitHub} \tabularnewline
\hline
แก่นแท้ & เครื่องมือควบคุมเวอร์ชัน (ซอฟต์แวร์) & แพลตฟอร์มจัดเก็บโค้ด (เว็บไซต์) \tabularnewline
ทำงานอยู่ที่ไหน & บนคอมพิวเตอร์ของตัวเอง & บนคลาวด์ (เข้าผ่านเบราว์เซอร์) \tabularnewline
หน้าที่หลัก & บันทึกการเปลี่ยนแปลงทุกครั้งของไฟล์ (แก้อย่างไร ใครแก้ แก้เมื่อไร) & เอารีโป Git ขึ้นคลาวด์ เก็บ โชว์ และทำงานร่วมกันได้สะดวก \tabularnewline
ต้องใช้เน็ตไหม & ไม่ต้อง ใช้แบบออฟไลน์ได้ & ต้องต่อเน็ตเพื่อเข้าใช้ \tabularnewline
เจ้าของ & ซอฟต์แวร์โอเพนซอร์ส สร้างโดย Linus Torvalds (พ่อของ Linux) & ผลิตภัณฑ์ในสังกัดไมโครซอฟต์ \tabularnewline
คู่แข่งสายเดียวกัน & SVN, Mercurial & GitLab, Gitee, Bitbucket \tabularnewline
\hline
\end{tabularx}
\end{center}

พูดสั้น ๆ:\textbf{Git ดูแล “การบันทึกการเปลี่ยนแปลงอย่างไร” ส่วน GitHub ดูแล “การเปลี่ยนแปลงเก็บไว้ที่ไหนและร่วมงานกันอย่างไร”}Git คือเครื่องมือ ทำงานอยู่บนเครื่องเรา GitHub คือแพลตฟอร์มที่เอารีโป Git ขึ้นคลาวด์ เพื่อความสะดวกในการร่วมมือและโชว์ผลงาน คุณใช้ Git อย่างเดียวโดยไม่ใช้ GitHub ก็ได้ (เช่น ทำเวอร์ชันเก็บโค้ดส่วนตัว) แต่ถ้าใช้ GitHub ต้องใช้ Git แน่นอน (เพราะ GitHub สร้างบนพื้นฐานของ Git นั่นเอง)

\booksection{สาม คอนเซปต์และฟีเจอร์หลักของ GitHub}

\booksubsection{คอนเซปต์พื้นฐาน}

ก่อนลงมือ มาทำความรู้จักคำที่ใช้บ่อยก่อน ไม่อย่างนั้นเปิดหน้า GitHub มาจะงงเป็นไข่แตก:

\begin{itemize}
  \item \textbf{รีโป (Repository หรือ Repo)}: หนึ่งโปรเจกต์ต่อหนึ่งรีโป ข้างในเก็บไฟล์โค้ดทั้งหมด ประวัติเวอร์ชัน และเอกสารของโปรเจกต์นั้น เข้าใจง่าย ๆ ว่าเป็น "โฟลเดอร์โปรเจกต์"
  \item \textbf{คอมมิต (Commit)}: บันทึกการแก้โค้ดหนึ่งครั้ง ทุกคอมมิตมีข้อความอธิบายกำกับว่า "รอบนี้แก้อะไร ทำไมถึงแก้" นี่คือรากฐานที่ทำให้ Git/GitHub "ย้อนเวลา" ได้
  \item \textbf{บรานช์ (Branch)}: เส้นทางพัฒนาที่แตกออกมาจากสายหลัก อย่างเช่นถ้าคุณอยากเพิ่มฟีเจอร์ใหม่ ก็เปิดบรานช์แยกไปพัฒนาคนละทาง พอเสร็จค่อยรวม (merge) กลับเข้าสายหลัก ไม่กระทบใคร
  \item \textbf{โคลน (Clone)}: ก๊อปรีโปบนคลาวด์มาไว้ที่คอมของเราแบบครบชุด จากนั้นก็แก้โค้ดบนเครื่องตัวเองได้เลย
  \item \textbf{พุช (Push) / พูล (Pull)}: Push คืออัปโหลดโค้ดจากเครื่องเราขึ้น GitHub ส่วน Pull คือดาวน์โหลดโค้ดล่าสุดจาก GitHub มาไว้ที่เครื่อง เป็นสองการเคลื่อนไหวหลักที่ทีมใช้ซิงก์โค้ดให้ตรงกัน
  \item \textbf{Pull Request (ย่อว่า PR)}: เมื่อคุณแก้โค้ดบนบรานช์เสร็จและอยากรวมเข้าบรานช์หลัก ก็ยื่น PR เพื่อให้ทีมรีวิวโค้ดก่อนแล้วค่อยรวม นี่คือกลไกหัวใจของการทำงานร่วมกันบน GitHub
  \item \textbf{อิชชู (Issue)}: ใช้จดบั๊ก ฟีเจอร์ที่ขอ และหัวข้ออภิปราย แต่ละ Issue เหมือนกระทู้สนทนาหนึ่งอัน มอบหมายผู้รับผิดชอบหรือติดป้ายกำกับได้
  \item \textbf{Fork}: ก๊อปรีโปของคนอื่นมาไว้ใต้บัญชีของตัวเอง จะแก้อะไรในสำเนาของตัวเองก็ได้อย่างอิสระ ไม่กระทบรีโปเดิม พอแก้เสร็จก็ส่ง PR กลับไปให้ผู้เขียนเดิม ขอให้รวมงานที่เราปรับปรุงเข้าไป
\end{itemize}

\booksubsection{ฟีเจอร์หลัก}

GitHub ให้มาไม่ใช่แค่ที่เก็บโค้ด แต่ยังรวมถึง:

\begin{itemize}
  \item \textbf{การควบคุมเวอร์ชัน}: บันทึกประวัติการแก้โค้ดทุกครั้ง พอแก้พังเมื่อไรก็ย้อนกลับไปเวอร์ชันปกติก่อนหน้าได้ทันที
  \item \textbf{การทำงานเป็นทีม}: รองรับหลายคนพัฒนาโปรเจกต์เดียวกันแบบขนานกันได้ หลายคนทำพร้อมกันก็ไม่ทับกัน ไม่ชนกัน
  \item \textbf{การติดตามปัญหา}: ใช้ Issue จดบั๊กและคำขอฟีเจอร์ ทำให้ทีมตามงานได้สะดวก
  \item \textbf{การจัดการโปรเจกต์}: มีเครื่องมือจัดการในตัว เช่น สิ่งที่ต้องทำ บอร์ดโปรเจกต์ (Projects) เหมือนมี Trello จิ๋วติดมาให้ด้วย
  \item \textbf{การทำงานอัตโนมัติ}: ผ่าน\textbf{GitHub Actions}ทำให้พอคอมมิตโค้ดแล้ว คอมไพล์ ทดสอบ แพ็ก และดีพลอยเองอัตโนมัติ เป็นเครื่องมือหัวใจของแนวคิด "CI/CD (การผสาน/การดีพลอยอย่างต่อเนื่อง)" ยุคใหม่
  \item \textbf{การจัดการเอกสาร}: README แนะนำโปรเจกต์ Wiki เก็บสถาปัตยกรรมละเอียดและข้อตกลงการพัฒนา แถม Pages ยังโฮสต์เว็บสถิตให้ได้เลย
  \item \textbf{การเก็บและโชว์โค้ด}: ทุกรีโปมีหน้าเว็บสวยงามให้ชม ไฮไลต์ไวยากรณ์โค้ดให้ด้วย เป็น "หน้าตา" ของผลงานสายเทคนิคของคุณ
  \item \textbf{ชุมชนโอเพนซอร์ส}: คุณกด Star (เก็บชอบ) โปรเจกต์ที่ชอบ Watch (ติดตาม) ข่าวคราวโปรเจกต์ Follow (ติดตาม) นักพัฒนา แล้วร่วมงานโอเพนซอร์สระดับโลกได้
\end{itemize}

\booksection{สี่ ลงทะเบียนและเริ่มใช้งาน: สร้างรีโปแรกของคุณใน 5 นาที}

\booksubsection{ลงทะเบียนบัญชี}

\begin{enumerate}
  \item เปิด\href{https://github.com}{github.com}แล้วกด "Sign up"
  \item กรอกอีเมล ตั้งรหัสผ่านและชื่อผู้ใช้ (ชื่อผู้ใช้จะโผล่ในที่อยู่โปรไฟล์ของคุณ แนะนำให้ใช้ภาษาอังกฤษ เช่น\texttt{github.\allowbreak com/\allowbreak yourname})
  \item ยืนยันอีเมลและผ่านการตรวจสอบว่าไม่ใช่บอท
  \item เลือกแพ็กฟรี (Free) ก็พอ สำหรับผู้ใช้ทั่วไปเกินพออยู่แล้ว
\end{enumerate}

\booksubsection{สร้างรีโปแรก}

\begin{enumerate}
  \item ล็อกอินแล้วกดปุ่ม "+" มุมขวาบน เลือก "New repository"
  \item ตั้งชื่อรีโป (Repository name) เช่น\texttt{my-\allowbreak first-\allowbreak project}
  \item เขียนคำอธิบายสั้น ๆ (Description) ให้คนอื่นรู้ว่ารีโปนี้ไว้ทำอะไร
  \item เลือกแบบสาธารณะ (Public) หรือส่วนตัว (Private) รีโปสาธารณะทุกคนเห็น ส่วนตัวเห็นเฉพาะตัวเองกับผู้ร่วมงาน บัญชีฟรีก็ทำรีโปส่วนตัวได้ไม่จำกัด
  \item ติ๊ก "Add a README file" — README คือ "คู่มือ" ของโปรเจกต์ แนะนำให้ทุกรีโปควรมี
  \item กด "Create repository" เป็นอันจบ รีโป GitHub แรกของคุณถือกำเนิดขึ้นแล้ว!
\end{enumerate}

\begin{tcbinfo}{รู้ไว้ใช่ว่า: README คืออะไร}
README.md คือไฟล์รูปแบบ Markdown ที่คนอื่นเห็นเป็นอย่างแรกเมื่อเปิดรีโปของคุณ README ที่ดีมักมี: แนะนำโปรเจกต์ วิธีติดตั้ง ตัวอย่างการใช้งาน สัญญาอนุญาต และช่องทางติดต่อ เขียน README ให้เก่งเป็นมารยาทพื้นฐานของโปรแกรมเมอร์ และเป็นประตูด่านแรกที่ตัดสินว่าคนอื่นจะใช้งานโปรเจกต์ของคุณได้หรือไม่

\end{tcbinfo}

\booksection{ห้า เวิร์กโฟลว์พื้นฐาน: จากเครื่องเราสู่คลาวด์}

หัวใจของเวิร์กโฟลว์ GitHub คือการจำวงจรนี้ให้ขึ้นใจ:\textbf{แก้ไฟล์ในเครื่อง สเตจ คอมมิต พุช คลาวด์อัปเดต}มาอธิบายแต่ละขั้นแบบภาษาคนทั่วไป:

\begin{enumerate}
  \item \textbf{แก้ไฟล์}: เปิดเอดิเตอร์บนคอมเขียนโค้ดหรือแก้เอกสาร
  \item \textbf{git add}: เพิ่มไฟล์ที่แก้เข้า "พื้นที่สเตจ" เทียบเท่ากับบอก Git ว่า "การแก้พวกนี้ เตรียมบันทึกไว้แล้ว"
  \item \textbf{git commit}: บันทึกการเปลี่ยนแปลงในพื้นที่สเตจอย่างเป็นทางการเป็น "คอมมิต" หนึ่งครั้ง พร้อมข้อความสั้น ๆ เช่น "แก้บั๊กหน้าล็อกอิน" ขั้นนี้เสร็จอยู่ในเครื่องเราล้วน ๆ
  \item \textbf{git push}: อัปโหลดคอมมิตจากเครื่องเราขึ้นคลาวด์ GitHub คนอื่นจึงจะเห็นงานได้
  \item \textbf{git pull}: ก่อนเริ่มทำงาน ดึงโค้ดล่าสุดจากคลาวด์มาก่อน เพื่อให้เครื่องเราตรงกับเพื่อนร่วมทีมเสมอ
\end{enumerate}

สำหรับมือใหม่ จำคำกลอนนี้ไว้:\textbf{“พูลก่อนแล้วค่อยแก้ แก้เสร็จก็แอด คอมมิตพร้อมข้อความ สุดท้ายพุชขึ้นคลาวด์”}ฝึกให้เป็นนิสัย ทำงานร่วมกับทีมก็จะไม่ปั่นป่วน

\booksection{หก เวิร์กโฟลว์ทีม: Pull Request ทำงานอย่างไร}

จุดแข็งที่สุดของ GitHub คือกลไกการทำงานร่วมกัน เวิร์กโฟลว์ทีมแบบคลาสสิกจะเป็นแบบนี้:

\begin{enumerate}
  \item \textbf{เปิดบรานช์}: แตกบรานช์ใหม่จากบรานช์หลัก (ปกติชื่อ main หรือ master) เช่น\texttt{feature/\allowbreak login}แล้วพัฒนาฟีเจอร์ใหม่บนบรานช์นี้
  \item \textbf{พัฒนาและคอมมิต}: เขียนโค้ดบนบรานช์ใหม่ คอมมิตหลายรอบเพื่อจดความคืบหน้า
  \item \textbf{พุชขึ้น GitHub}: พุชบรานช์ขึ้นรีโปบนคลาวด์
  \item \textbf{ยื่น Pull Request}: เปิด PR บนหน้าเว็บ GitHub ขอรวมบรานช์ของคุณเข้าบรานช์หลัก เขียนให้ชัดว่า PR นี้แก้อะไรและทำไมถึงแก้
  \item \textbf{รีวิวโค้ด (Code Review)}: เพื่อนร่วมทีมอ่านโค้ดของคุณ เสนอความเห็นหรือแนะนำให้แก้ คุณคอมมิตต่อตามฟีดแบ็กได้ PR จะอัปเดตเองอัตโนมัติ
  \item \textbf{รวมโค้ด (Merge)}: รีวิวผ่านแล้ว กด "Merge" รวมโค้ดเข้าบรานช์หลัก ผลงานของคุณเข้าโปรเจกต์อย่างเป็นทางการแล้ว!
\end{enumerate}

เวิร์กโฟลว์นี้รับประกันว่า\textbf{โค้ดทุกชิ้นผ่านการรีวิวก่อนเข้าบรานช์หลักเสมอ}ความเสี่ยงพลาดจึงต่ำลงมาก โปรเจกต์โอเพนซอร์สก็ใช้กลไกนี้รับงานจากคนแปลกหน้าทั่วโลก — คุณยื่น PR เข้าโปรเจกต์โอเพนซอร์สดัง ๆ ก็ได้ พอถูกรับ คุณก็กลายเป็นผู้ร่วมพัฒนาโปรเจกต์นั้นแล้ว

\booksection{เจ็ด สถานการณ์ใช้งาน GitHub แบบคลาสสิก}

\booksubsection{นักศึกษากับการแข่งขันและโปรเจกต์รายวิชา}

นักศึกษาคอมพิวเตอร์ที่ลงแข่ง ACM แฮ็กะธอน หรือทำโปรเจกต์จบ แทบทุกคนต้องพึ่ง GitHub มันทำหน้าที่ทั้งแพลตฟอร์มทำงานร่วมกันของทีม และหน้าต่างโชว์ผลงาน — กรรมการแข่งเปิดดูรีโปของคุณได้ตรง ๆ เห็นทั้งคุณภาพโค้ดและกระบวนการทำงานร่วมกัน หลายบริษัทรับสมัครงานก็ขอลิงก์โปรไฟล์ GitHub มาเป็นหลักฐานฝีมือทางเทคนิคด้วย

\booksubsection{คลังความรู้ส่วนตัวและที่เก็บโน้ต}

คนที่ไม่ใช่โปรแกรมเมอร์ก็ใช้ GitHub ได้ หลายคนใช้เก็บโน้ตเรียน บล็อกสายเทคนิค (คู่กับ GitHub Pages ทำเว็บได้ฟรี) หรืออีบุ๊ก เขียนด้วย Markdown + จัดเวอร์ชันด้วย Git + สำรองบนคลาวด์ ครบสามอย่างในที่เดียว

\booksubsection{ร่วมพัฒนาโปรเจกต์โอเพนซอร์ส}

GitHub คือฐานทัพหลักของกระแสโอเพนซอร์สทั่วโลก คุณ Fork โปรเจกต์ที่ชอบ แก้แล้วส่ง PR กลับคืนชุมชน หรือจะแค่ Star เก็บโปรเจกต์ดี ๆ ไว้ตามอัปเดตก็ได้ การร่วมโอเพนซอร์สเป็นทางที่ดีที่สุดทั้งฝึกฝีมือและสะสมชื่อเสียง

\booksubsection{ดีพลอยอัตโนมัติกับ CI/CD}

ด้วยพลังของ GitHub Actions ทำ “พุชโค้ดปุ๊บ ทดสอบเอง + ดีพลอยขึ้นเซิร์ฟเวอร์เอง” แบบอัตโนมัติเต็มระบบได้ บล็อกส่วนตัวหรือเว็บเอกสารจำนวนมากตั้งมาแบบนี้ เขียนบทความเสร็จพุชหนึ่งที เว็บก็อัปเดตเองทันที

\booksubsection{สำรองโค้ดและซิงก์หลายเครื่อง}

ใช้รีโปโค้ดเป็นคลาวด์ไดรฟ์สิ เปลี่ยนคอม ทำงานหลายเครื่อง เมื่อไรก็ git pull หนึ่งครั้งได้โค้ดล่าสุด ไม่ต้องแบกแฟลชไดรฟ์ก๊อปไปก๊อปมาอีกต่อไป

\booksection{แปด คำแนะนำสำหรับมือใหม่}

\begin{itemize}
  \item \textbf{ข้อความคอมมิตต้องเขียนให้ชัด}: อย่าเขียนขี้เรื่องแบบ "update" commit message ที่ดีคือแบบ "แก้ปัญหาผู้ใช้ล็อกอินแล้วเด้งไปหน้าผิด" ผ่านไปครึ่งปีกลับมาอ่านก็ยังรู้ว่าวันนั้นทำอะไรไว้
  \item \textbf{คอมมิตบ่อย ๆ ก้าวเล็กแต่ไว}: ทำฟีเจอร์เล็ก ๆ เสร็จชิ้นหนึ่งก็คอมมิตหนึ่งครั้ง ไม่ใช่เขียนทั้งวันค่อยมาคอมมิตรอบเดียว แบบนี้พอมีปัญหาจะไล่หาต้นตอได้ง่ายกว่า
  \item \textbf{พัฒนาบนบรานช์ อย่าแก้บรานช์หลักตรง ๆ}: แม้จะเป็นโปรเจกต์ส่วนตัว ก็ฝึกเปิดบรานช์ไว้ ต่อให้ต้องทำงานเป็นทีมในอนาคตก็จะไม่งงมืองงเท้า
  \item \textbf{ใช้ .gitignore ให้เป็น}: ไฟล์นี้บอก Git ว่าไฟล์ไหนไม่ต้องเอาเข้าการจัดเวอร์ชัน (เช่น ไฟล์ตั้งค่าที่มีรหัสผ่าน ผลลัพธ์คอมไพล์ โฟลเดอร์ไลบรารี) กันไม่ให้เอาขยะขึ้นไปกองไว้
  \item \textbf{ชมโปรเจกต์โอเพนซอร์สชั้นเลิศให้มากขึ้น}: แวะไปเล่น GitHub Trending ดูโครงสร้างโค้ด วิธีเขียน README และเวิร์กโฟลว์ของโปรเจกต์ดาวเยอะ ๆ เป็นวิธีฟรีที่พัฒนาฝีมือเร็วที่สุดอย่างหนึ่ง
  \item \textbf{ขอ GitHub Student Pack}: ยืนยันสถานะนักศึกษาแล้วจะได้เครดิตใช้เครื่องมือและบริการเสียเงินมากมายฟรี (Copilot โดเมน คลาวด์ ฯลฯ) มูลค่าหลายพันดอลลาร์
\end{itemize}

\begin{tcbwarn}{คำเตือนด้านความปลอดภัย}
\begin{itemize}
  \item ห้ามคอมมิตรหัสผ่าน คีย์ API สตริงเชื่อมต่อฐานข้อมูล หรือข้อมูลละเอียดอ่อนอื่นขึ้น GitHub เด็ดขาด ยิ่งเป็นรีโปสาธารณะยิ่งอันตราย — พอรั่วไหล เหล่าบอทไต่เว็บจะกวาดไปภายในเสี้ยววินาที ผลที่ตามมาร้ายแรงมาก
  \item ถ้าพลาดคอมมิตข้อมูลละเอียดอ่อนไปแล้ว การลบไฟล์อย่างเดียวไม่พอ (ประวัติยังเก็บมันอยู่) ต้องใช้\texttt{git filter-\allowbreak branch}หรือเครื่องมือเฉพาะทางล้างประวัติ และรีบเปลี่ยนคีย์ที่รั่วทันที
  \item รีโปส่วนตัวก็ไม่ได้ปลอดภัยแบบสัมบูรณ์ ข้อมูลละเอียดอ่อนควรจัดการด้วยตัวแปรสภาพแวดล้อม (environment variables) ไม่เอาเข้ารีโปโค้ดเลยดีที่สุด
\end{itemize}

\end{tcbwarn}

\begin{tcbinfo}{ข้อควรรู้เรื่องการเข้าถึง GitHub จากในจีน}
การเข้าถึง GitHub จากในจีนบางช่วงก็ไม่เสถียร อาจเจอเปิดไม่ได้หรือโคลนช้า ลองทำดู: (1) ใช้มิเรอร์สายในประเทศ (เช่น GitClone, ghproxy) เร่งความเร็วการโคลน (2) ตั้งค่า Hosts หรือใช้พร็อกซี (3) ถ้างานหลักเน้นทำงานร่วมกันในประเทศ ก็พิจารณาตัวแทนสายจีนอย่าง\textbf{Gitee (Mayun)}ฟีเจอร์ใกล้เคียงกันและเข้าถึงไวกว่า หลายทีมโฮสต์ทั้ง GitHub และ Gitee พร้อมกัน ได้ทั้งความสากลและความเร็วในประเทศ

\end{tcbinfo}

\textbf{เทคนิคขั้นสูง: สร้างบล็อกส่วนตัวฟรีด้วย GitHub Pages}\par

GitHub Pages ให้คุณโฮสต์เว็บสถิตได้ฟรี คู่กับตัวสร้างเว็บสถิตอย่าง Jekyll, Hexo, Hugo ก็ตั้งบล็อกส่วนตัวได้ง่าย ๆ เขียนบทความ Markdown เสร็จพุชเข้ารีโป บล็อกก็อัปเดตเอง แถมมี HTTPS และรองรับโดเมนของตัวเองด้วย บล็อกของบล็อกเกอร์สายเทคนิคนับไม่ถ้วนตั้งมาแบบนี้ ไม่มีค่าใช้จ่าย ไม่มีโฆษณา ควบคุมเวอร์ชันได้

\booksection{เก้า สรุป: ทำไมทุกคนควรรู้จัก GitHub}

GitHub ไม่ใช่แค่ “เครื่องมือของโปรแกรมเมอร์” อีกต่อไป มันมากกว่านั้น เป็น\textbf{วิถีการทำงานแบบยุคดิจิทัล}: การจัดเวอร์ชันทำให้ทุกอย่างย้อนดูได้ กลไกการทำงานร่วมกันทำให้ทีมวิ่งได้เต็มสปีด จิตวิญญาณโอเพนซอร์สทำให้ความรู้ไหลได้อย่างอิสระ แม้คุณไม่เขียนโค้ด การเข้าใจคอนเซปต์อย่าง “คอมมิต บรานช์ รวมโค้ด” ก็ช่วยให้คุณจัดการทุกสิ่งที่ต้อง “แก้ซ้ำ ๆ + คนหลายคนร่วมทำ” ได้ดีขึ้น — ตั้งแต่วิทยานิพนธ์ เอกสาร โน้ต ไปจนถึงแผนงาน

สำหรับนักศึกษาคอมพิวเตอร์ GitHub ยิ่งเป็นทักษะต้องมี: มันคือพอร์ตโฟลิโอ แพลตฟอร์มทำงานร่วมกัน คลังแหล่งเรียนรู้ และ “เรซูเม่” ที่หนักแน่นที่สุดตอนหางานในอนาคต ลงทะเบียนบัญชีให้เร็ว ค่อย ๆ อัปงานรายวิชาและโปรเจกต์เล็ก ๆ ของตัวเองขึ้นไปสะสมทีละนิด อีกไม่กี่ปีคุณจะมีประวัติทางเทคนิคที่ทำให้คนสัมภาษณ์ตาลุกวาว

\begin{tcbpractice}{มาลงมือกัน: เปิดประสบการณ์แรกกับ GitHub}
\begin{enumerate}
  \item ลงทะเบียนบัญชี GitHub แล้วตกแต่งโปรไฟล์ให้ครบ (รูปหน้าปก แนะนำตัว สายเทคนิคที่สนใจ)
  \item สร้างรีโปสาธารณะชื่อ\texttt{hello-\allowbreak github}พร้อมติ๊ก "Add a README file"
  \item เข้าไปในรีโป กดปุ่มแก้ไขของ README.md (ไอคอนดินสอ) เขียนแนะนำตัวสั้น ๆ ด้วย Markdown แล้วคอมมิตการเปลี่ยนแปลง
  \item ค้นหาคำที่สนใจในแถบค้นหาของ GitHub (เช่น\texttt{python-\allowbreak tutorial}) แล้ว Star โปรเจกต์ที่ชอบสักสองอัน
  \item (ขั้นสูง) ติดตั้ง Git บนคอม โคลนรีโปมาไว้ที่เครื่อง แก้ README แล้วพุชกลับขึ้นคลาวด์ ให้ครบหนึ่งวงจรเวิร์กโฟลว์ "เครื่องเรา คลาวด์"
\end{enumerate}

\end{tcbpractice}
